% ============================================================
%  Section 1, Segment 1C — Hybrid Picking / 混合拨弦
%  原书页码 18-21
% ============================================================
\chapter{Hybrid Picking / 混合拨弦}
\label{ch:hybrid}

\eng{Hybrid picking is where we make use of a pick and fingers at the same time. It is not as
prominent in math rock as other techniques we have explored but still very useful. I for one
love this technique and employ it regularly. I like to use it because I can get the best of both
worlds; the strengths of fingerpicking and having the option of using a pick at the same time. I
implore you to learn this technique if you have not been doing so already.}

\chn{混合拨弦是当我们同时使用拨片和手指的技巧。它在数学摇滚中的地位不如前面探索的技巧那样显著，但仍然非常实用。就我而言，我很喜欢这项技巧，并经常使用。它让我能够兼得两者优势：指弹的灵活性和拨片的力度。我强烈建议你学习这项技术，即使你还没有开始使用它。}

\eng{This technique, when developed, can help you play fast precise licks, especially ones involving
string skipping with a lot less effort. In addition, you can combine the pick with fingers for
plucking chords ( just like fingerpicking). Overall, it is a great compromise between using only
fingers or only a pick.}

\chn{当这项技巧熟练后，它可以帮助你演奏快速而精准的乐句，尤其是那些涉及跳弦的乐句，而且所需的努力更少。此外，你还可以将拨片与手指结合起来拨奏和弦（就像指弹一样）。总体而言，它是在只用手指和只用拨片之间的一个很好的折中方案。}

\paragraph{Examples of guitarists who use hybrid picking}
\begin{itemize}[leftmargin=*]
  \item \eng{Mario \& Erick of CHON – Book}
  \chn{CHON 的 Mario 与 Erick — 《Book》}
  \item \eng{Tim \& Scott of Polyphia – Look but Don't Touch Playthrough}
  \chn{Polyphia 的 Tim 与 Scott — 《Look but Don't Touch》演奏版}
  \item \eng{Alan of Tangled Hair – Yeah… It does look like a spider}
  \chn{Tangled Hair 的 Alan — 《Yeah… It does look like a spider》}
\end{itemize}

\eng{Pay close attention to the singing parts and how all the strings resonate simultaneously when Alan plucks them, in contrast to when he strums during the louder sections. Additionally, starting from the two-minute mark, observe how Alan's use of hybrid picking enables him to execute that section, achieving the distinct sound it produces due to the accessibility afforded by this technique.}

\chn{特别注意演奏中的连奏部分，以及当 Alan 拨奏时所有琴弦同时共鸣的效果，这与他在更响亮部分进行扫弦时形成了对比。另外，从两分钟处开始，观察 Alan 如何利用混合拨弦来演奏该部分，使其产生这种特殊音色，这正是该技巧所带来的便利性。}

\paragraph{Hybrid Picking Dexterity Exercises}
\eng{Below you will find a series of exercises I have designed to help you become accustomed with the technique.}

\chn{下面你会找到一系列我设计的练习，帮助你熟悉这项技巧。}

\eng{To clarify, you will be holding a pick between your thumb and index finger (as you normally would) and using your remaining free fingers (middle, ring, and little finger) to pick notes. Compared to fingerpicking, I have noticed I rarely use my little finger when hybrid picking. That is not to say you will not need to use it, just a reminder that it is fine not to.}

\chn{说明一下，你将把拨片夹在拇指和食指之间（像平常那样），并使用其余空闲的手指（中指、无名指和小指）拨奏音符。与指弹相比，我注意到混合拨弦时我很少使用小指。这并不是说你绝对不用它，只是提醒你不使用它也没问题。}

\eng{To begin with, exercise 3.A is designed to get you started with the string skipping potentials of hybrid picking. I have labelled which fingers to use (remember Mac for this one). The lowest string will be played with your pick.}

\chn{首先，练习 3.A 的设计是让你入门混合拨弦的跳弦潜力。我已标注出应使用的手指（记住 Mac 的指法）。最低音弦由拨片弹奏。}

\eng{Keep your picking hand steady by resting it on the bridge, and notice how you hand and wrist are not jumping about like you would if you were to pick all these notes using only a pick.}

\chn{将拨弦手稳稳放在琴桥上，观察你的手和手腕不会像仅用拨片弹奏时那样乱动。}

\eng{I designed Exercise 3.B to get your fingers familiar with picking up and down in an arch like motion. This kind of arching melody motion is common in math rock.}

\chn{练习 3.B 的设计是让你的手指熟悉类似拱形的上下拨动。这种拱形旋律运动在数学摇滚中很常见。}

\eng{Bonus tip for exercise 3.B: once you are comfortable with playing it precisely, try picking up and down as fast as you can (without hurting yourself of course!). You can develop this into a quick raking like action across the strings (like sweep picking but with fingers included). It is very satisfying once you can do it!}

\chn{练习 3.B 的额外提示：一旦你能精准演奏，尝试尽可能快地上下拨动（当然不要伤到自己！）。你可以把它发展成一种快速扫弦般的动作（类似扫弦，但加入了手指）。一旦你能做到，会非常有成就感！}

\eng{Exercise 3.C is a slight variation on the previous exercise. It is designed to get your fingers moving in different patterns.}

\chn{练习 3.C 是前一个练习的一个小变化。它的目的是让你的手指以不同模式移动。}

\eng{Exercise 3.D applies a picking pattern along with some arpeggios.}

\chn{练习 3.D 将拨弦模式与琶音结合起来。}

\eng{Let's move on to plucking strings at the same time. Exercise 3.E gets you started with plucking two strings at a time using different finger combinations along with a pick. This may feel a little strange at first but that is completely normal.}

\chn{接下来，我们开始同时拨奏多根琴弦。练习 3.E 帮助你开始练习使用不同手指组合配合拨片同时拨奏两根弦。一开始可能会有点奇怪，但这是完全正常的。}

\eng{Exercise 3.F expands on the previous exercise by including the use of two fingers with a pick and in the last bar the inclusion of all three spare fingers!}

\chn{练习 3.F 在前一个练习的基础上扩展到使用两个手指配合拨片，在最后一小节中加入了全部三个空闲手指！}

\eng{Exercise 3.G uses just two notes from chord shapes (highest and lowest in this case) to create some interesting ideas.}

\chn{练习 3.G 只使用和弦形状中的两个音（在这个例子中为最高音和最低音），以创造一些有趣的想法。}

\paragraph{Test Yourself: Hybrid Picking Exercises in Context}
\eng{Exercise 3.H is an adaptation of the main chord progression to the song English Man in New York by Sting. Use your pick to play the lowest note of each chord and use your fingers to pluck the rest of the chord shape.}

\chn{练习 3.H 是 Sting 的《English Man in New York》主和弦进行的改编。用拨片弹每个和弦的最低音，用手指拨奏剩余的和弦形状。}

\eng{Lastly, exercise 3.I uses a repeating ascending pattern in triplets. Again, like exercises before, use the pick for the lower notes. Notice in the last bar, that by using hybrid picking, the string skipping becomes somewhat easier than trying to play it with just a pick.}

\chn{最后，练习 3.I 使用三连音的重复上行模式。同前几个练习一样，用拨片弹低音。注意在最后一个小节，通过混合拨弦，跳弦变得比仅用拨片演奏更容易。}
